%% abstract.tex

\begin{abstract}
Large language model (LLM) agents are increasingly deployed with \emph{self-healing} mechanisms---automated recovery pipelines that capture execution failures and re-inject diagnostic context into subsequent agent iterations. While prior work has extensively studied prompt injection against the forward execution path, the \emph{recovery channel} has remained an overlooked attack surface. Diagnostic logs and semantic core dumps are routinely granted implicit system-level trust, and the recovery actions they trigger bypass the permission checks enforced on the forward path. We introduce \textbf{Recovery-Channel Injection (RCI)}, a class of attacks in which adversarial content originating from external sources traverses the recovery pipeline and is re-stamped with system-level trust signals.

We employ a two-stage methodology common in security research: a \emph{white-box testbed} for mechanism attribution and a \emph{black-box production framework} for external validity. Because production-grade agent frameworks (LangGraph, AutoGen, SWE-agent) do not expose sufficient internal hooks for fine-grained mechanism attribution---e.g., isolating whether content-level provenance or structural privilege enforcement is driving a given defense outcome---we construct \system{}, a multi-agent operating system prototype with a multi-tier recovery architecture, as a white-box testbed. On \system{} we isolate and study two structural attack instances: (i)~\emph{trust escalation} via high-trust template re-injection (\code{ReflectionInjectionRecovery}), where externally-originated error traces inherit a default \code{SYSTEM\_GENERATED} label; and (ii)~\emph{privilege escalation} via zero-permission topology mutation (\code{TopologyMutationStrategy}), where a poisoned core dump induces an attacker-suggested role that reaches the node-replacement primitive without permission re-check. We propose a lightweight, OS-principled defense combining \trustorigin{} Tagging (four-level content provenance) with a unified \reauthgate{} enforcing \emph{Privilege Downgrade on Recovery}. In a real-LLM evaluation on the testbed (400 API calls to DeepSeek-V4-Flash and Qwen2.5-72B-Instruct; temperature~0.7), the defense reduces V1 ASR from \textbf{0.57 to 0.07} (Wilson 95\% CI $[0.47, 0.66] \rightarrow [0.03, 0.14]$; $p < 10^{-6}$, Fisher's exact) and V2 ASR from \textbf{1.00 to 0.00} (deterministic: the privilege-downgrade invariant rejects out-of-set roles by construction), with \textbf{0.0} FAR on benign failures and sub-millisecond overhead. A stronger-model evaluation on \textbf{gpt-oss-120b} (open-weights, not closed-weight frontier) reveals that content-level V1 defense degrades on stronger instruction-following models (defended ASR: 0.07 $\rightarrow$ 0.50), while structural V2 defense remains at 0.00---suggesting structural enforcement becomes \emph{necessary} as models improve.

The testbed results establish the \emph{mechanism}; external validity comes from LangGraph, a production framework not developed by the authors. A forward-path experiment on \code{create\_react\_agent} using entirely library-default components (no author-constructed code) achieves 0.88 V1 pooled ASR ($[0.76, 0.94]$), confirming that production frameworks lack provenance tagging on the tool-output path---though this is the classic indirect injection vector, not RCI per se. A supplementary \emph{real agentic-loop} experiment (where the LLM actually executes the canary tool rather than merely describing intent) confirms ecological validity: baseline ASR rises to \textbf{0.96} ($[0.80, 0.99]$), validating that the proxy measurement used in the main evaluation is conservative. A recovery-path experiment on \code{StateGraph} using an author-constructed recovery node (following common Reflexion-style patterns) achieves V1 pooled ASR of \textbf{0.76} ($[0.57, 0.89]$), reduced to \textbf{0.04} ($[0.01, 0.20]$) by the untrusted-framing defense, demonstrating that the RCI mechanism is realizable on a third-party API; V2 achieves 0.54 pooled ASR ($[0.40, 0.67]$), reduced to 0.00 by the \reauthgate{} analog. A \emph{critical experiment} establishes RCI as an independent attack surface: with three forward-path defenses fully enabled (input sanitizer + LLM injection detector + tool-permission check), the RCI attack still achieves 0.68 ASR ($[0.54, 0.79]$), statistically indistinguishable from the undefended 0.74 ($p=0.660$)---forward-path defenses operate on user input, while the adversarial payload enters via the tool's failure, which is never inspected. To address the concern that the LangGraph recovery node is author-constructed, we additionally conducted a \emph{source-level wild-case study} of four externally-maintained, production-deployed agent codebases (Reflexion, MetaGPT, Aider, OpenHands; combined $\sim$105k GitHub stars): three of four ship retry logic with the structural signature of the V1 trust-escalation trampoline (no provenance tag $\wedge$ high-trust frame). For Reflexion, we further \emph{confirmed the vulnerability end-to-end}: an exploit against Reflexion's actual \code{generic\_generate\_func\_impl} retry path (commit \texttt{218cf0e})---using its verbatim system prompt, few-shot example, and \code{parse\_code\_block}---achieves 0.52 pooled ASR ($[0.39, 0.65]$, $p = 3.5 \times 10^{-10}$ vs.\ 0\% FPR control), with two of five variants at 100\% ASR. The evidence chain is: (i) forward-path (library-default components lack provenance tagging), (ii) recovery-path (RCI realizable on third-party APIs), (iii) forward-defense-bypass (RCI succeeds under full forward-path defenses), (iv) wild-case (structural provenance gap in shipped production code) + confirmed end-to-end exploit on Reflexion's actual code, with the testbed providing mechanism attribution. The burden of proof rests on the two pieces of evidence that use \emph{no author-constructed code}---the library-default LangGraph forward-path experiment (0.88 ASR) and the confirmed Reflexion end-to-end exploit (0.52 ASR on real commit \texttt{218cf0e})---with \system{} serving purely as a mechanism-attribution auxiliary, not as evidence that production systems are vulnerable. A supplementary evaluation on a closed-weight model (accessed via the official API) confirms that the V1 content-level defense degrades: defended ASR is 0.36 (vs.\ 0.07 on DeepSeek/Qwen), and the baseline-vs-defended difference is \emph{statistically significant} ($p=0.027$, $N=50$ per config)---the defense reduces V1 ASR but the residual 0.36 means more than one in three attacks still succeeds; the V2 structural defense is LLM-independent by construction.
\end{abstract}
